% Paste into Overleaf in place of the current Premise section. Requires amsmath.
% Matches aircraft.py as built. Items marked TODO need your group's decision.

\section{The Premise}

\subsection{The Question}
An airliner flying a fixed route wants to burn as little fuel as possible. Two choices made
early in the flight dominate the answer: how steeply it climbs (flight-path angle $\gamma$)
and how high it cruises ($h_{\text{cruise}}$). A steep climb reaches thin air quickly, where
parasitic drag is low, but demands high thrust. A shallow climb is gentle but spends a long time
at low altitude, where drag is high. A high cruise altitude lowers parasitic drag but raises the
lift coefficient and so the induced drag, and the Mach limit caps the speed. These effects compete, so we
expect an optimum.

\textbf{Question.} For flights from Chicago O'Hare (ORD) to New York JFK ($1{,}191$~km),
Copenhagen ($6{,}873$~km) and Hong Kong ($12{,}543$~km), which climb angle $\gamma$ and cruise
altitude $h_{\text{cruise}}$ minimize the total fuel burned, and how does the optimum shift with
trip length?

\subsection{The Model}
The aircraft is a point mass in a vertical plane with state $\mathbf{y}=[V,\,h,\,x,\,m]^\top$
(airspeed, altitude, ground distance, mass). It starts at $h=0$ and flies two phases:
a \textbf{climb} at constant $\gamma$ until $h=h_{\text{cruise}}$, then a level \textbf{cruise}
($\gamma=0$) until $x=X$. Descent and landing are neglected, and every route starts at the same takeoff
mass $m_0$.

\subsubsection{Atmosphere}
We use the NASA Glenn two-layer model ($h$ in m, $T$ in $^\circ$C, $p$ in kPa),
\begin{align}
  h \le 11000: \quad & T = 15.04 - 0.00649\,h, \quad
     p = 101.29\left(\tfrac{T+273.1}{288.08}\right)^{5.256}, \\
  h > 11000: \quad & T = T_{11},\quad p = p_{11}\,e^{-0.000157\,(h-11000)}, \qquad
     \rho = \frac{p}{0.2869\,(T+273.1)},
\end{align}
where $T_{11}$ and $p_{11}$ are the troposphere values at $11$~km. This differs from the published
fit by under $0.2\%$ and makes $T$, $p$ and $\rho$ continuous at $11$~km. (The published fit has a $0.1\%$ density jump, which limits
the order of any integrator for trips crossing it.)

\subsubsection{Aerodynamics, Speed Schedule and Thrust}
\begin{align}
  C_L &= \frac{2\,m\,g_0\cos\gamma}{\rho V^2 S}, \qquad
  C_D = C_{D0} + k\,C_L^2, \qquad
  D = \tfrac12\rho V^2 S\,C_D .
\end{align}
The pilot flies the speed that minimizes drag at the current altitude and mass, limited by Mach $M_{\max}$:
\begin{equation}
  V_{\text{tgt}}(h,m) = \min\!\left\{\sqrt{\frac{2mg_0}{\rho S}}\left(\frac{k}{C_{D0}}\right)^{1/4},\; M_{\max}\,a(h)\right\}.
\end{equation}
Thrust holds $\gamma$ and tracks this speed with a proportional speed-hold law, and fuel flow follows thrust:
\begin{equation}
  T = D + m g_0\sin\gamma + m K\,(V_{\text{tgt}} - V), \qquad \dot m_{\text{fuel}} = \frac{T}{I_{sp}\,g_0}.
\end{equation}

\subsubsection{The Initial Value Problem}
\begin{align}
  \dot V &= \frac{T-D}{m} - g_0\sin\gamma, &
  \dot h &= V\sin\gamma, &
  \dot x &= V\cos\gamma, &
  \dot m &= -\frac{T}{I_{sp}\,g_0},
\end{align}
with $\mathbf{y}(0)=[V_{\text{tgt}}(0,m_0),\,0,\,0,\,m_0]^\top$.
The fuel burned is $m_0 - m(t_{\text{end}})$.

\subsubsection{Feasibility Constraints (masking)}
A $(\gamma,h_{\text{cruise}})$ pair is \emph{infeasible} if, at any time in the flight, the required thrust exceeds $T_{\max}$,
the lift coefficient exceeds $C_{L,\max}$, or the fuel runs out ($m \le m_{\text{dry}}$). Infeasible cases are masked out of
the results.

\subsection{Why Numerical Methods Are Needed}
In the climb, $\rho$ depends on $h(t)$, and $C_L$, $D$ and $T$ depend nonlinearly on the evolving $V$, $h$ and $m$,
so the system has no closed-form solution. Only special cases do (level cruise at fixed speed, and a linear
constant-density limit); we use those as independent checks on the code. Answering the premise needs hundreds of simulations
(one per $(\gamma,h_{\text{cruise}},X)$ combination), so accuracy per unit cost matters.

\subsection{Parameters}
\begin{itemize}
  \item \textbf{Varied:} $\gamma$, $h_{\text{cruise}}$, trip distance $X$.
  \item \textbf{Fixed (illustrative widebody, not calibrated to any real aircraft):}
        $m_0 = 400$~t, $m_{\text{dry}} = 210$~t, $S = 550$~m$^2$, $C_{D0}=0.020$, $k=0.045$, $I_{sp}=6000$~s,
        $T_{\max}=600$~kN, $C_{L,\max}=1.2$, $M_{\max}=0.85$, $K = 0.05$~s$^{-1}$.
  \item \textbf{Output:} total fuel burned $\Delta m_{\text{fuel}}(\gamma,h_{\text{cruise}};X)$.
\end{itemize}
